\section{The dual fibration and spectral Nyquist}

This section records claim nodes and dependency edges. Paper-grade development in
\texttt{2026-08-26-paper-outline-two-fibrations-nyquist.md}. Entry fact, recorded as
prose because it is the section's reason to exist: sinc reconstruction is non-monotone
in the pointwise order (negative lobes), so it is provably not a leg of any position-side
Galois connection --- and it is exactly the fundamental-domain transversal of the
\emph{dual} fibration below. The box and the sinc are one construction on the two sides
of Fourier duality.

\begin{proposition}[The dual fibration]\label{thm:dual-fibration}\lean{Obsmetrologia.Valuation.DualFibration}\notready\uses{thm:sampling-fibration, thm:finiteFibration}
\emph{(Open frontier --- stated, not formalized.)} For a finite abelian carrier $G$ with
sampling subgroup $H \le G$: the quotient of the dual carrier by the annihilator,
$q : \hat{G} \to \hat{G}/H^{\perp}$, with $H^{\perp}$-cosets as listed fibers --- a
\cref{thm:finiteFibration} on the nose, Fourier-dual to the position-side
\cref{thm:sampling-fibration}. First instance: $G = \texttt{Fin}\,n$,
$H$ the multiples of $d \mid n$, coset space concretely $\texttt{Fin}\,d$; at the
possibilistic layer no characters are needed, only \texttt{Fin} arithmetic, and the
Crystals stratum's finite cyclic period is the natural carrier to reuse.
\end{proposition}

\begin{theorem}[Aliasing is fiber collision]\label{thm:aliasing-fiber-collision}\lean{Obsmetrologia.Valuation.sample_spectrum_existsF, Obsmetrologia.Valuation.aliasing_iff_fiber_collision}\notready\uses{thm:dual-fibration, thm:existsF-adjunction}
\emph{(Open frontier --- stated, not formalized.)} Support-level Poisson summation:
sampling on $H$ acts on spectral support as $\exists$ along $q$ --- the sampled
spectrum's support is the fiber-join of the original's. Aliasing is a fiber containing
two or more support points: the lossy left adjoint collapsing distinct fiber elements,
the same $\exists$ of \cref{thm:existsF-adjunction} met on the dual side.
\end{theorem}

\begin{theorem}[Spectral Nyquist, possibilistic layer]\label{thm:spectral-nyquist}\lean{Obsmetrologia.Valuation.recoverable_iff_fiber_injective, Obsmetrologia.Valuation.filter_is_transversal}\notready\uses{thm:aliasing-fiber-collision, thm:pixelize-adjoint-triple}
\emph{(Open frontier --- stated, not formalized.)} A signal is exactly recoverable from
its $H$-samples iff its spectral support is fiber-injective under $q$ --- a partial
section of the quotient. Reconstruction filters correspond to transversals of $q$; the
``ideal'' filter is the fundamental-domain transversal, whose spectral action is
meet-with-the-band (an interior operator) and whose position-side kernel is the
Dirichlet kernel, the finite sinc. Completes \cref{thm:pixelize-adjoint-triple}'s
reading of Smith: the little square and the sinc are the fundamental-domain indicators
of the two sides, exchanged by the duality. The amplitude layer (aliasing \emph{sums}
components) is explicitly gated on a value category with addition, KCBS-style:
possibilistic layer first.
\end{theorem}

\begin{theorem}[Conflation swaps the failure modes]\label{thm:conflation-swaps-failure-modes}\lean{Obsmetrologia.Valuation.duality_carries_kbot_to_ktop}\notready\uses{thm:spectral-nyquist, thm:sampling-failure-bilattice, thm:conflate-orders, thm:klein-four-symmetry}
\emph{(Open frontier --- stated, not formalized.)} The aliasing event reads as
$\bot_k$-drift on the position side (under-determination: many signals, one sample set)
and as $\top_k$-drift on the frequency side (superposition: two components, one fiber
--- told-both). The involution carrying the one failure reading to the other is
conflation, the second generator of \cref{thm:klein-four-symmetry}'s $K_{4}$ ---
pinning both Klein-four generators to concrete sampling phenomena. (The companion
conjecture --- the lens transform as the first generator on ray space --- is held in
the note's discussion section, dervish discipline: no node until the involutions are
constructed and checked against \cref{thm:conflate-orders}.)
\end{theorem}

\begin{theorem}[Projection-slice is Beck--Chevalley]\label{thm:fourier-slice-beck-chevalley}\lean{Obsmetrologia.Valuation.projection_slice_beck_chevalley}\notready\uses{thm:beck-chevalley, thm:dual-fibration, thm:lenticular-fiber-sample}
\emph{(Open frontier --- stated, not formalized.)} Ray space $\texttt{Domain} \times D$
carries three canonical aggregations --- pixel ($\exists$ over direction at fixed
position), hogel (the direction fiber at one position,
\cref{thm:lenticular-fiber-sample}), sinogram ($\exists$ over position at fixed
direction, the Radon projection row) --- as the legs of one two-axis fibration. The
projection-slice theorem (marginalize-then-transform equals transform-then-slice;
classically Radon, photographically Ng 2005) is the base-change square of the two
fibrations: reindexing along the duality commuting with the fiber quantifiers, an
instance of \cref{thm:beck-chevalley}. Potentially the cheapest headline theorem in the
chapter: the square is already green; the work is instantiating it.
\end{theorem}
